\section{Введение}
Семантическая сегментация сопоставляет каждому пикселю изображения класс объекта: дорога, автомобиль, пешеход и другие категории. Пиксельная разметка реальных сцен требует значительных трудозатрат. Синтетические сцены позволяют получать изображения и соответствующие метки из виртуальной среды, однако различия текстур, освещения, геометрии и частот классов затрудняют перенос модели на фотографии.

Доменная адаптация без целевой разметки (unsupervised domain adaptation, UDA) использует размеченные изображения исходного домена и неразмеченные изображения целевого домена. Важен не только выбор сегментационной сети, но и критерий, позволяющий извлечь обучающий сигнал из целевых данных. В AdaptSegNet используется состязательное согласование карт предсказаний~\cite{tsai}; DACS сочетает псевдоразметку с междоменным смешиванием изображений~\cite{dacs}; DAFormer показывает существенное влияние архитектуры и стратегии обучения~\cite{daformer}.

Итоговые метрики этих работ не позволяют ранжировать отдельные математические принципы: одновременно меняются модель, начальные веса, аугментации и правила обучения. Цель работы --- сопоставить механизмы адаптации и сформулировать условия их проверки. Аналитический результат включает систематизацию критериев, их ограничений и правил сравнения опубликованных оценок.
